\chapter{Core Concepts}
\label{ch:core-concepts}

\newthought{You have already met all five of them.} If you ran the quick start, you
created a \emph{project}, wrote paths in a \emph{token language}, had a sampler propose
points in \emph{u-space}, watched twenty \emph{Samples} go through a lifecycle, and read
back a table of \emph{observables}. What you did not have was the vocabulary. This
chapter supplies it --- once, precisely, for the whole book --- and after it every
\textsc{yaml} key you meet and every file a run writes is recognisable as one of five
things.

\section{One point, end to end}
\label{sec:one-point}

Before the definitions, the whole machine in one pass. Take a single point out of the
quick-start scan and follow it from the moment the sampler invents it to the moment you
open it in a spreadsheet. Table~\ref{tab:one-point} follows it stage by stage.

\begin{table}[ht]
  \footnotesize
  \begingroup
  \setlength{\tabcolsep}{0pt}
  \begin{tabular}{@{}p{0.255\linewidth}p{0.213\linewidth}p{0.532\linewidth}@{}}
    \toprule
    \textbf{Stage} & \textbf{It is now} & \textbf{Where it lives} \\
    \midrule
    sampler proposes &
      \code{u = (0.37,\allowbreak{} 0.82)} &
      a task on the Redis queue: \textsc{uuid}, u-coords, execution plan \\
    Worker maps it &
      \code{x = 0.37,\allowbreak{} y = 0.82} &
      the Worker's memory: physical values your workflow can use \\
    workflow runs &
      \code{+ z = 57.24} &
      the observables dictionary, one entry richer \\
    likelihood scores &
      \code{+ LogL\_Z,\allowbreak{} LogL} &
      the same dictionary; \code{LogL} is what the sampler steers by \\
    Archiver files it &
      one row &
      \file{DATABASE/\allowbreak{}samples.hdf5}, exported to \file{samples.csv} \\
    \bottomrule
  \end{tabular}
  \endgroup
  \caption[One point, from proposal to spreadsheet.]{One point, from proposal to
  spreadsheet. Each row is one of this chapter's five ideas at work.}
  \label{tab:one-point}
\end{table}

Two things in that table are worth pausing on, because they are the two most common
sources of confusion later.

First, \textbf{the point exists twice}: once as \code{u}, the sampler's currency, and
once as \code{x}, your physics. In the quick start those two happened to be numerically
identical --- both variables were \yamlkey{Flat} on $[0,1]$, so the map between them was
the identity, which is exactly why you did not notice there was a map at all. Change one
line, to \yamlkey{Log} on $[100, 10000]$, and \code{u = 0.37} becomes \code{x = 549.54}.
Section~\ref{sec:uspace} is about that map and about which of your \textsc{yaml} numbers
live on which side of it.

Second, \textbf{the dictionary only ever grows}. It starts as the parameters, gains one
entry per calculator or opera output, gains the likelihood terms last, and is then
written out whole. Nothing edits it in place and nothing hands values around behind its
back. Section~\ref{sec:observables} makes that precise, and it is the reason expressions
anywhere in a card can refer to anything computed before them.

\section{Projects and the project root}
\label{sec:projects}

The first question a run must answer is \emph{where am I}, and it answers it before
anything else happens. Every path in your card, and every file the run writes, hangs off
one directory: the \emph{project root}.

Normal usage runs from a standalone project directory, not from a checkout of the
framework. A project is any directory holding your task cards, data, external tools, and
outputs. When \cli{Jarvis} loads a task card, it resolves the root in this order, and
stops at the first hit:

\begin{enumerate}
  \item The environment variables \code{JARVIS\_HEP\_TASK\_ROOT} or
        \code{JHEP\_TASK\_ROOT}, when set. This one wins over everything, which makes it
        the right lever on a batch system where the card is not where you think it is.
  \item Otherwise, walk up from the task card's directory looking for a marker file:
        \file{.jarvis-project.json} or \file{jarvis.project.yaml}.
  \item Failing that, a parent directory containing \file{bin/}; failing that, the task
        card's own directory.
\end{enumerate}

That last clause is why the quick start worked with nothing but a loose \textsc{yaml}
file: no marker anywhere, so the card's own directory became the root, and \file{outputs/}
appeared beside it. \cli{Jarvis project create} scaffolds a properly marked project
instead (Chapter~\ref{ch:project-cli}), and the difference shows the moment you move the
card: a marked project keeps its root, a loose card takes its root with it.

Everything a run writes is project-local. The output root is
\file{<project-root>/outputs/<scan-name>/}, where the scan name resolves as
\yamlkey{Scan.name} $\rightarrow$ \yamlkey{scan\_name} $\rightarrow$ \code{default}, and
can be overridden wholesale with the top-level key \yamlkey{task\_result\_dir}
(Chapter~\ref{ch:task-yaml}).

\section{Path markers and tokens}
\label{sec:tokens}

Task cards never need absolute paths --- a card that hard-codes \file{/home/you/} is a
card nobody else can run, and one you cannot move to a cluster. A small token language
keeps every card portable.

The important structure is that tokens resolve in two \emph{phases}, at two different
times, in two different processes. Knowing which phase a token belongs to tells you
exactly where it may appear.

\textbf{Phase 1 --- static, resolved once at load time} in the control process. These are
the answers that are the same for the whole run: \marker{\&J/} for the project root,
\code{\$\{Scan:name\}} for the scan, \code{\$\{LibDeps:<name>\}} for where an external
tool was installed, and the bare name of any tool you registered.

\textbf{Phase 2 --- per-Sample, resolved inside the Worker} just before execution. These
are the three things that cannot be known until a point exists: \token{@SampleID} (its
\textsc{uuid}), \token{@Sdir} (its own directory), and \token{@PackID} (the calculator
slot the Worker holds).

\subsection{Watching one string get resolved}

A single string can carry both phases, and each half is filled in by whichever process
knows the answer. Written in the card:

\begin{terminal}[as written in the task card]
\begin{jcmd}
${LibDeps:SPheno}/bin/SPheno @Sdir/LesHouches.in.@SampleID
\end{jcmd}
\end{terminal}

The control process resolves the static half once, at load time, for every Sample that
will ever run. The \token{@Sdir} and \token{@SampleID} survive untouched --- they cannot
be known yet, because no Sample exists. Later, in a Worker, moments before the program
starts:

\begin{terminal}[what the Worker actually ran]
\begin{jcmd}
/scratch/tools/SPheno-4.0.5/bin/SPheno \
  /work/proj/outputs/scan/SAMPLE/000003/8f2c.../LesHouches.in.8f2c...
\end{jcmd}
\end{terminal}

This is why the phases are worth remembering rather than looking up: a Phase-1 token is
resolved once and costs nothing per Sample, while a Phase-2 token is what makes two
Samples running the same program at the same instant not collide. Putting a token in the
wrong place does not silently produce a literal string --- Phase 2 refuses:

\begin{terminal}
\begin{jout}
Phase 2 cannot resolve static tokens in stage 'execution' field 'command': ...
\end{jout}
\end{terminal}

That is the whole idea, and it is enough to read the task-card chapters with.
Chapter~\ref{ch:command-system} gives the complete token tables, the rules about where
each one may appear, and how registered executables work.

\section{u-space and physical space}
\label{sec:uspace}

Samplers do not think in your physics units, and this is deliberate. Every sampler
proposes points in \emph{u-space} --- the unit hypercube $[0,1]^d$, one dimension per
entry of \yamlkey{Sampling.Variables} --- and a \emph{mapper} converts each coordinate to
the physical value your workflow sees, using that variable's declared distribution
(Chapter~\ref{ch:variables}). Figure~\ref{fig:uspace} is the whole arrangement in three
steps.

\begin{marginfigure}
\centering
\smallcaps{u-space} $[0,1]^d$\\[2pt]
$\big\downarrow$ \emph{mapper (per variable)}\\[2pt]
\smallcaps{physical point} $x$\\[2pt]
$\big\downarrow$ \emph{workflow}\\[2pt]
\smallcaps{observables}
\caption{Samplers propose in the unit cube; the Worker maps to physical values.}
\label{fig:uspace}
\end{marginfigure}

The same \code{u} means completely different physics depending on what you declared ---
compare the two right-hand columns of Table~\ref{tab:uspace-mapping}:

\begin{table}[ht]
  \footnotesize
  \begingroup
  \setlength{\tabcolsep}{0pt}
  \begin{tabular}{@{}p{0.255\linewidth}p{0.213\linewidth}p{0.532\linewidth}@{}}
    \toprule
    \textbf{u} & \textbf{Flat 0--5} & \textbf{Log 100--10000} \\
    \midrule
    0.00 & 0.00 & 100.00 \\
    0.25 & 1.25 & 316.23 \\
    0.37 & 1.85 & 549.54 \\
    0.50 & 2.50 & 1000.00 \\
    1.00 & 5.00 & 10000.00 \\
    \bottomrule
  \end{tabular}
  \endgroup
  \caption[One u-coordinate, two declared distributions.]{One u-coordinate, two declared
  distributions. Half of u-space is below 1000 in the log case and below 2.5 in the flat
  one --- and the sampler cannot tell the difference.}
  \label{tab:uspace-mapping}
\end{table}

\begin{jwarn}
\textbf{Sampler geometry keys are u-space quantities, never physical ones.}
\sampler{Bridson}'s \yamlkey{Bounds.Radius} and \sampler{AdaptiveBridson}'s
\yamlkey{initial\_radius} and \yamlkey{min\_radius} are fractions of the unit cube. Writing
\yamlkey{Radius: 50} for a mass you think of in \textsc{gev} does not request 50\,GeV of
spacing --- it requests fifty times the entire width of every axis, and you get one
point.
\end{jwarn}

The split also pays for itself in the distributed runtime: the only payload a task
carries through Redis is its u-coordinates. The u\,$\rightarrow$\,x map is rebuilt
identically inside every Worker, so proposing stays cheap and the wire stays light
(Chapter~\ref{ch:distributed-runtime}).

The exception is the \sampler{CSV} method, which replays already-physical points: there
the mapper is a no-op and parameters travel with the task
(Chapter~\ref{ch:csv-sampler}).

\section{The Sample lifecycle}
\label{sec:sample-lifecycle}

A \emph{Sample} is one parameter point together with everything it accumulates: a
\textsc{uuid}, u-coordinates, an execution plan, a working directory, a private log, and
finally a result record. Its life crosses three processes, and each one leaves something
on disk you can go and look at:

\begin{enumerate}
  \item \textbf{Proposed} (Sampler). The sampler mints a deterministic \textsc{uuid} and
        pushes a light task dictionary --- \textsc{uuid}, u-coordinates, execution plan
        --- onto the Redis queue. \emph{Leaves behind:} a queue entry, and a line in
        \file{sampler.log}.
  \item \textbf{Executed} (one Worker). A Worker pulls the task, maps
        u\,$\rightarrow$\,x, allocates the Sample a slot in the current \emph{bucket}
        (below), and runs the execution plan layer by layer: same-layer calculators run
        concurrently, layers run in order, operas run in-process, and the likelihood is
        evaluated last. Success or failure, the Worker submits a result record and moves
        on. \emph{Leaves behind:} \file{SAMPLE/<bucket>/<uuid>/} with the Sample's own
        artifacts and its private log.
  \item \textbf{Archived} (Archiver). The Archiver batches finished Samples into the
        \textsc{hdf5} database and packs completed buckets. \emph{Leaves behind:} a row in
        \file{DATABASE/\allowbreak{}samples.hdf5}, and eventually
        \file{<bucket>.\allowbreak{}tar.gz}.
\end{enumerate}

\begin{jnote}
Failure is a first-class outcome, not an exception path: a failed Sample keeps its
working directory and its detailed log for forensics, is counted in the statistics, and
is archived as a failed record. One bad point never aborts the scan --- which is the
whole reason a scan can run overnight unattended.
\end{jnote}

\paragraph{Buckets.} Per-sample artifacts live under
\file{SAMPLE/<bucket>/<uuid>/}, where buckets are numbered directories
(\file{000001}, \file{000002}, \ldots) holding a bounded number of Samples --- 200 by
default. A full, fully-archived bucket is packed to \file{<bucket>.tar.gz}. The reason is
mundane and important: a directory with a million entries makes a network filesystem
crawl, and \code{ls} in it stops returning. The knobs live under
\yamlkey{EnvReqs.V2.sample\_directory} (Chapter~\ref{ch:envreqs}).

\section{Observables}
\label{sec:observables}

As a Sample executes, it accumulates a single flat dictionary of named values, the
\emph{observables}. For the quick-start point of Section~\ref{sec:one-point}, the finished
dictionary is exactly this:

\begin{terminal}[one Sample's observables]
\begin{jout}
x        = 0.37        # parameter, from the mapper
y        = 0.82        # parameter, from the mapper
z        = 57.2436     # opera output: (sin(x)*cos(y) + 2)^5
LogL_Z   = 57.2436     # named likelihood term, expression: z
LogL     = 57.2436     # the total, written for you
\end{jout}
\end{terminal}

Five names, and every one of them is addressable by name from anywhere else in the card.
That is what makes this dictionary the lingua franca of the whole system:

\begin{itemize}
  \item calculator \yamlkey{output} specs and opera \yamlkey{output} lists \emph{add}
        entries to it (Chapters~\ref{ch:calculators} and~\ref{ch:operas});
  \item every expression --- opera inputs, \yamlkey{Dump} variables,
        \yamlkey{LogLikelihood} terms, sampler \yamlkey{selection} cuts,
        \sampler{AdaptiveBridson} targets --- is evaluated \emph{over} it
        (Chapter~\ref{ch:expressions});
  \item one record per Sample lands in \file{DATABASE/samples.hdf5}, and
        \file{samples.csv} exports every column (Chapter~\ref{ch:outputs}).
\end{itemize}

Note \code{LogL} in that listing: you named one term, \code{LogL\_Z}, and the runtime
also wrote the total. \code{LogL} is what the sampler steers by, and it is the sum of
every term you declared --- with a single term the two coincide, which is why the quick
start's numbers look redundant.

Name collisions are resolved by order: later writers win. Keep observable names unique
and expression-friendly --- they double as symbols.

\section{The frozen contracts}
\label{sec:contracts}

A few behaviors are \emph{contracts}: stable surfaces you may build analysis pipelines
on, revised only with a major release. If you are writing a script that will still be run
next year, these are the parts to write it against.

\begin{itemize}
  \item The task-card schema --- every key documented in Part~\ref{part:task-card}; new
        keys are always optional.
  \item The \cli{Jarvis} command surface and its exit codes (Chapter~\ref{ch:cli}).
  \item The output layout: \file{DATABASE/samples.hdf5} record structure, the
        full-column \file{samples.csv}, the \file{SAMPLE/} bucket layout, and
        \file{run\_summary} fields.
  \item Checkpoint and resume behavior (Chapter~\ref{ch:checkpoint}).
\end{itemize}

\begin{jtip}
Every task card in Part~\ref{part:task-card} is the same four sentences in \textsc{yaml}:
\emph{name a project-local scan} (\S\ref{sec:projects}), \emph{declare variables, which
defines a u-space} (\S\ref{sec:uspace}), \emph{declare a workflow, which fills the
observables} (\S\ref{sec:observables}), and \emph{score them}. The runtime does the rest
--- moving Samples through their lifecycle (\S\ref{sec:sample-lifecycle}) along paths
that stayed portable (\S\ref{sec:tokens}). When a key in a later chapter seems to come
from nowhere, ask which of those four sentences it belongs to.
\end{jtip}
